\section{Задание}
\begin{enumerate}
    \item[1)] Для данной табличной функции вычислить приближённое значение в точке интерполяции с помощью многочлена Лагранжа и схемы Эйткена.

    \item[2)] Для табличной функции из задания 1 вычислить приближённые значения в точках интерполяции с помощью подходящего многочлена Ньютона с разделёнными разностями.

    \item[3)] Для данной табличной функции вычислить приближённые значения в точках интерполяции с помощью подходящего многочлена Ньютона с конечными разностями (точки $x^{(1)}$, $x^{(2)}$) и наиболее подходящего из многочленов Гаусса, Стирлинга или Бесселя (точка $x$).
\end{enumerate}

\section{Цель работы}
Научиться приближённо вычислять табличную функцию с помощью интерполяционных многочленов Лагранжа, Ньютона, Гаусса, Стирлинга и Бесселя.

\section{Ход работы}

\textit{Перед выполнением работы отметим, что все исходные данные взяты из варианта 47.}

\subsection{Задание 1}
Требуется вычислить приближённое значение функции в данной точке $x = 4.96$ с помощью
\begin{enumerate}
    \item[а)] многочлена Лагранжа,
    \item[б)] схемы Эйткена,
\end{enumerate}
используя таблицу значений \ref{table:values}.

\begin{table}[H]
    \centering
    \caption{Значения $x_i,y_i$ для построения интерполяционного многочлена}
    \label{table:values}
    \begin{tblr}{
        colspec = {X[-1,l]X[c]X[c]X[c]X[c]X[c]|},
        hlines,
        vlines,
        rows = {abovesep=8pt,belowsep=8pt, $},
        column{1} = {leftsep=8pt},
        }
        x_i & 1.41    & 2.31   & 4.13   & 5.31   & 6.01   \\
        y_i & -1.4156 & 2.3901 & 3.0567 & 0.9812 & 2.7569 \\
    \end{tblr}
\end{table}

\subsubsection*{а) с помощью многочлена Лагранжа}

\noindentФормула Лагранжа для системы $n$ узлов:
\[
L_n(x)=\frac{(x- x_1) \dots (x- x_n)}{(x_0 - x_1) \dots (x_0 - x_n)} y_0+ \dots + \frac{(x- x_0) \dots (x- x_n)}{(x_i - x_0) \dots (x_i - x_n)} y_i~+
\]
\[
+ \dots + \frac{(x- x_0) \dots (x- x_{n-1})}{(x_n - x_0) \dots (x_n - x_{n-1})} y_n.
\]

\noindentСокращённая форма формулы Лагранжа для системы $n$ узлов:
\[
L_n(x)= \sum_{i=0}^{n} y_i l_{n,i}(x), \quad \text{где } l_{n,i}(x)= \prod_{\substack{j=0 \\ j \neq i}}^{n} \frac{x- x_j}{x_i - x_j}.
\]

\begin{listing}[H]
\caption{Вычисление интерполяционного многочлена Лагранжа в точке 4.96}
\vspace{-20pt}
\label{lst:task1_lagrange}
\begin{minted}[linenos, numbersep=10pt, firstnumber=1]{matlab}
x = [1.41, 2.31, 4.13, 5.31, 6.01];
y = [-1.4156, 2.3901, 3.0567, 0.9812, 2.7569];

x_interp = 4.96;
L = 0;

n = length(x);
for i = 1:n
    L_i = 1;
    for j = 1:n;
        if i ~= j
            L_i = L_i * (x_interp - x(j)) / (x(i) - x(j))
        end
    end
    L = L + y(i) * L_i
end


fprintf('%f\n', L)
\end{minted}
\end{listing}

В результате работы программы было полученно приблизательное значение функции в точке $x = 4.96$.

\[
L_6(4.96) = 1.289912 \approx 1.29.
\]

\noindentОтвет: $f(4.96) = 1.29.$

\subsubsection*{б) с помощью схемы Эйткена}

Cхема Эйткена предназначена для пошагового рекуррентного вычисления значения интерполяционного многочлена Лагранжа $L_n$ в точке $x$.

\noindentНачальные условия для реккурентной схемы:

\[
L^i_0(x) = y_i, i = 0, 1, \dots, n.
\]

\noindentФормула для первого шага:
\[
L_{1}^{i,i+1}(x)= \frac{1}{x_{i+1} - x_i} \begin{vmatrix} L_0^i(x) & x_i - x \\ L_0^{i+1}(x) & x_{i+1} - x \end{vmatrix} = \frac{1}{x_{i+1} - x_i} \begin{vmatrix} y_i & x_i - x \\ y_{i+1} & x_{i+1} - x \end{vmatrix},
\]
где $i= 0, 1,\dots, n- 1$.

\noindentФормула для $k$-того шага:
\[
L_k^{i,i+1,\dots,i+k}(x)= \frac{1}{x_{i+k} - x_i} \begin{vmatrix} L_{k-1}^{i,\dots,i+k-1}(x) & x_i - x \\ L_{k-1}^{i,\dots,i+k}(x) & x_{i+k} - x \end{vmatrix},
\]
где $i= 0, 1,\dots, n- k$; $k= 1,\dots, n$.

\noindentФормула для $n$-того шага:
\[
L_n(x) = L_n^{0,\dots,n}(x)= \frac{1}{x_{n} - x_0} \begin{vmatrix} L_{n-1}^{0,\dots,n-1}(x) & x_0 - x \\ L_{n-1}^{0,\dots,n}(x) & x_{n} - x \end{vmatrix}.
\]

На листинге \ref{lst:task2} приведён код программы, вычисляющей приближённое значение функции в точке $x=4.96$ методом Эйткена.

\begin{listing}[H]
\caption{Реализация вычислительной схемы Эйткена в MATLAB}
\vspace{-20pt}
\label{lst:task2}
\begin{minted}[linenos, numbersep=10pt, firstnumber=1]{matlab}
x = [1.41, 2.31, 4.13, 5.31, 6.01];
y = [-1.4156, 2.3901, 3.0567, 0.9812, 2.7569];

x_interp = 4.96;
n = length(x);

L_table = zeros(n, n);
L_table(:, 1) = y';

for k = 2:n
    for i = 1:(n-k+1)
         M = [ L_table(i, k-1),     x(i) - x_interp;
              L_table(i+1, k-1),   x(i+k-1) - x_interp ];

        L_table(i, k) = det(M) / (x(i+k-1) - x(i));
    end
end

result = L_table(1, n);
fprintf('%f\n', result);
\end{minted}
\end{listing}

В результате работы программы было полученно приблизательное значение функции в точке $x = 4.96$.

\[
L_6(4.96) = 1.289912 \approx 1.29.
\]

\noindentОтвет: $f(4.96) = 1.29.$

\subsection{Задание 2}
Требуется вычислить значение функции в точках $x^{(1)}= 6.17$ и $x^{(2)}= 2.11$, с помощью одного из многочленов Ньютона с разделёнными разностями, используя таблицу значений \ref{table:values}.

Пусть функция $y=f(x)$ задана таблицей с попарно несовпадающими узлами, которые пронумерованы в произвольном порядке. Для такой таблицы можно вычислить разделенные разности. Разделенная разность первого порядка функции $f$ в узлах $x_i, x_{i+1}$, вычисляется по формуле

\[
f(x;x_{i+1}) = \frac{y_{i+1}-y_i}{x_{i+1}-x_i} = \frac{f(x_{i+1})-f(x_i)}{x_{i+1}-x_i},~i = 0, \dots, n-1.
\]

Разделенная разность $k$-го порядка в узлах $x_i, x_{i+1}, \dots, x_{i+k}$ вычисляется по формуле

\[
f(x;\dots;x_{i+k}) = \frac{f(x_{i+1}; \dots; x_{i+k})-f(x_i; \dots; x_{i+k-1})}{x_{i+k}-x_i},~i = 0, \dots, n-k.
\]

Интерполяционным многочленом Ньютона с разделенными разностями первого вида (многочленом интерполяции вперёд) называется многочлен

\[
P_n^{(1)}(x) = f(x_0) + f(x_0;x_1)(x - x_0) + f(x_0;x_1;x_2)(x - x_0)(x - x_1)~+
\]

\[
+ \dots + f(x_0;\dots;x_n)(x - x_0)\dots(x - x_{n-1}) =
\]

\[
= f(x_0) + \sum_{i=1}^n f(x_0;\dots;x_i)(x - x_0)\dots(x - x_{i-1}).
\]

Применяется для вычисления значения заданной функции в точке, близкой к началу таблицы.

Интерполяционным многочленом Ньютона с разделенными разностями второго вида (многочленом интерполяции назад) называется многочлен

\[
P_n^{(2)}(x) = f(x_n) + f(x_{n-1};x_n)(x - x_n) + f (x_{n-2};x_{n-1};x_n)(x - x_n)(x - x_{n-1})~+
\]

\[
+ \dots + f(x_0;\dots;x_n)(x - x_n)\dots(x - x_1) =
\]

\[
= f(x_n) + \sum_{i=1}^n f(x_{n-i};\dots;x_n)(x - x_n)\dots(x - x_{n-i+1}).
\]

Применяется для вычисления значения заданной функции в точке, близкой к концу таблицы.

Составим таблицу разделённых разностей (табл. \ref{table:diff}).

\begin{table}[H]
    \centering
    \caption{Таблица разделённых разностей}
    \label{table:diff}
    \begin{tblr}{
        colspec = {X[c]X[c]X[c]X[c]X[c]X[c]},
        hlines,
        vlines,
        rows = {abovesep=5pt,belowsep=5pt, $},
        cell{1}{2-3} = {font=\small},
        cell{1}{4-6} = {font=\footnotesize}
    }
        x_i &
        f(x_i) &
        f(x_i;x_{i+1}) &
        f(x_i;\dots;x_{i+2}) &
        f(x_i;\dots;x_{i+3}) &
        f(x_0;\dots;x_4) \\

        1.4100 & -1.4156 &         &         &         &        \\
               &         &  4.2286 &         &         &        \\
        2.3100 &  2.3901 &         & -1.4200 &         &        \\
               &         &  0.3663 &         & 0.1825  &        \\
        4.1300 &  3.0567 &         & -0.7084 &         & 0.1362 \\
               &         & -1.7589 &         & 0.8090  &        \\
        5.3100 &  0.9812 &         &  2.2849 &         &        \\
               &         &  2.5367 &         &         &        \\
        6.0100 &  2.7569 &         &         &         &        \\
    \end{tblr}
\end{table}

На листинге \ref{lst:task2_t} приведена программа, вычисляющая данную таблицу.

\begin{listing}[H]
\caption{Вычисление таблицы разделённых разностей}
\vspace{-20pt}
\label{lst:task2_t}
\begin{minted}[linenos, numbersep=10pt, firstnumber=1]{matlab}
x = [1.41, 2.31, 4.13, 5.31, 6.01];
y = [-1.4156, 2.3901, 3.0567, 0.9812, 2.7569];

function D = calcDivDiffTable(x, y)
    n = length(x);
    D = zeros(n, n+1);
    D(:, 1) = x';
    D(:, 2) = y';

    for j = 3:n+1
        for i = 1:n-j+2
            D(i, j) = (D(i+1, j-1) - D(i, j-1)) / (D(i+j-2, 1) - D(i, 1));
        end
    end
end

D = calcDivDiffTable(x, y);

disp('Divided Difference Table:');
disp(D);
\end{minted}
\end{listing}

Рассмотрим указанные в задании точки $x^{(1)} = 6.17$ и $x^{(2)} = 2.11$. Точка $x^{(1)}$ находится близко к концу таблицы, следовательно, для вычисления значения функции в ней следует использовать многочлен Ньютона второго вида.

На листинге \ref{lst:task2_1} приведён код программы, вычисляющий значение функции в точке $x^{(1)} = 6.17$ с помощью интерполяционного многочлена Ньютона второго вида.

\begin{listing}[H]
\caption{Вычисление значения функции в точке $x^{(1)} = 6.17$ помощью многочлена Ньютона второго вида}
\vspace{-20pt}
\label{lst:task2_1}
\begin{minted}[linenos, numbersep=10pt, firstnumber=1]{matlab}
x = [1.41, 2.31, 4.13, 5.31, 6.01];
y = [-1.4156, 2.3901, 3.0567, 0.9812, 2.7569];

D = calcDivDiffTable(x, y);
n = length(x);
x_interp = 6.17;

P = D(n, 2);
term = 1;

for i = 1:n-1
    term = term * (x_interp - x(n-i+1));
    P = P + D(n-i, i+2) * term;
end

fprintf('%f\n', P)
\end{minted}
\end{listing}

Таким образом,

\[
P_6^{(1)}(x^{(1)}) = P_6^{(1)}(6.17) = 3.851845 \approx 3.85.
\]

Точка $x^{(2)} = 2.11$, напротив, находится близко к началу таблицы, следовательно, для вычисления значения функции в ней следует использовать многочлен Ньютона первого вида.

На листинге \ref{lst:task2_2} приведён код программы, вычисляющий значение функции в точке $x^{(2)}$ с помощью интерполяционного многочлена Ньютона первого вида.

\begin{listing}[H]
\caption{Вычисление значения функции в точке $x^{(2)} = 2.11$ с помощью многочлена Ньютона первого вида}
\vspace{-20pt}
\label{lst:task2_2}
\begin{minted}[linenos, numbersep=10pt, firstnumber=1]{matlab}
x = [1.41, 2.31, 4.13, 5.31, 6.01];
y = [-1.4156, 2.3901, 3.0567, 0.9812, 2.7569];

D = calcDivDiffTable(x, y);
n = length(x);
x_interp = 2.11;

P = D(1, 2);
term = 1;

for i = 1:n-1
    term = term * (x_interp - x(i));
    P = P + D(1, i+2) * term;
end

fprintf('%f\n', P)
\end{minted}
\end{listing}

Таким образом,

\[
P_6^{(2)}(x^{(2)}) = P_6^{(2)}(2.11) = 1.671522 \approx 1.67.
\]

\subsection{Задание 3}

Для заданной табличной функции $f(x) = \arccos(\log_2(\frac{x}{2}))$ требуется вычислить значение функции в точках интерполяции $x^{(1)}= 1.32$ и $x^{(2)}= 3.84$, с помощью подходяшего многочленоа Ньютона с конечными разностями, и наиболее подходящего
из многочленов Гаусса, Стирлинга или Бесселя в точке $x = 2.2$.

Соотавим таблицу \ref{table:values_2} значений функции на отрезке $[1;4]$ с шагом $h = 0.5$.


\begin{table}[H]
    \centering
    \caption{Значения функции $f$ на отрезке $[1;4]$ с шаром $h = 0.5$}
    \label{table:values_2}
    \begin{tblr}{
        colspec = {X[-1,l]X[c]X[c]X[c]X[c]X[c]X[c]X[c]},
        hlines,
        vlines,
        rows = {abovesep=8pt,belowsep=8pt, $},
        column{1} = {leftsep=8pt},
        }
        x_i & 1.0000 & 1.5000 & 2.0000 & 2.5000 & 3.0000 & 3.5000 & 4.0000 \\
        y_i & 3.1416 & 1.9988 & 1.5708 & 1.2430 & 0.9460 & 0.6311 &      0 \\
    \end{tblr}
\end{table}

На листинге \ref{lst:task3_values} приведена пограмма, вычисляющая таблицу знчений функции $f$.

\begin{listing}[H]
\caption{Вычисление таблицы значений для функции $f(x) = \arccos(\log_2(\frac{x}{2}))$}
\vspace{-20pt}
\label{lst:task3_values}
\begin{minted}[linenos, numbersep=10pt, firstnumber=1]{matlab}
nodes = 7;

V = zeros(2, nodes);
h = 0.5;

for i = 1:nodes
    x = 1 + (i - 1) * h;
    V(1, i) = x;
    V(2, i) = acos(log2(x / 2));
end

first = V(1, 1);
last = V(1, end);

disp(V);
\end{minted}
\end{listing}


Пусть функция $y = f(x)$ задана таблицей, в которой узлы идут в возрастающем порядке с постоянным шагом $h = x_i - x_{i-1}$ ($x_i = x_0 + ih, i = 0, 1, \dots, n.$). Тогда для $f$ можно вычислить конечные разности.

\noindentКонечная разность функции $f$ на $i$-ом узле:

\[
\Delta y_i = y_{i+1} - y_i = f(x_{i+1}) - f(x_i), i = 0, \dots, n-1.
\]

\noindentРазности порядков выше первого определяются рекуррентно через предыдущие.

\[
\Delta^k y_i = \Delta^{k-1} y_{i+1} - \Delta^{k-1} y_i, i = 0, \dots, n-k; k = 2, \dots, n.
\]

\noindentИнтерполяционный многочлен Ньютона с конечными разностями первого вида:

\[
P_n^{(1)}(x) = P_n^{(1)}(x_0 + qh) = f(x_0) + \frac{\Delta y_0}{1!}q + \dots + \frac{\Delta^n y_0}{n!}q(q-1)\dots(q-n+1) =
\]

\[
= f(x_0) + \sum_{i = 1}^{n} \frac{\Delta^i y_0}{i!}q(q-1)\dots(q-i+1).
\]

\[
q = \frac{x-x_0}{h}.
\]

\noindentИнтерполяционный многочлен Ньютона с конечными разностями второго вида:

\[
P_n^{(2)}(x) = P_n^{(2)}(x_n + qh) = f(x_n) + \frac{\Delta y_{n-1}}{1!}q + \dots + \frac{\Delta^n y_0}{n!}q(q+1)\dots(q+n-1) =
\]

\[
= f(x_n) + \sum_{i = 1}^{n} \frac{\Delta^i y_{n-i}}{i!}q(q+1)\dots(q+i-1).
\]

\[
q = \frac{x-x_n}{h}.
\]

Вычислим таблицу \ref{table:diff2} конечных разностей.
\begin{table}[H]
    \centering
    \caption{Таблица конечных разностей}
    \label{table:diff2}
    \begin{tblr}{
        colspec = {X[c]X[c]X[c]X[c]X[c]X[c]X[c]X[c]},
        hlines,
        vlines,
        rows = {abovesep=5pt,belowsep=5pt, font=\small, $},
        row{1} = {font=\normal}
    }
        x_i &
        y_i &
        \Delta y_i &
        \Delta^2 y_i &
        \Delta^3 y_i &
        \Delta^4 y_i &
        \Delta^5 y_i &
        \Delta^6 y_0 \\

        1.0000 &  3.1416 &         &         &         &         &         &        \\
               &         & -1.1428 &         &         &         &         &        \\
        1.5000 &  1.9988 &         &  0.7148 &         &         &         &        \\
               &         & -0.4280 &         & -0.6146 &         &         &        \\
        2.0000 &  1.5708 &         &  0.1002 &         &  0.5451 &         &        \\
               &         & -0.3278 &         & -0.0695 &         & -0.5240 &        \\
        2.5000 &  1.2430 &         &  0.0307 &         &  0.0211 &         & 0.2528 \\
               &         & -0.2971 &         & -0.0485 &         & -0.2712 &        \\
        3.0000 &  0.9460 &         & -0.0178 &         & -0.2501 &         &        \\
               &         & -0.3148 &         & -0.2986 &         &         &        \\
        3.5000 &  0.6311 &         & -0.3163 &         &         &         &        \\
               &         & -0.6311 &         &         &         &         &        \\
        4.0000 &       0 &         &         &         &         &         &        \\
    \end{tblr}
\end{table}

На листинге \ref{lst:task3_t} приведена пограмма, вычисляющая данную таблицу.

\begin{listing}[H]
\caption{Вычисление таблицы конечных разностей на основе таблицы значений для функции $f(x) = \arccos(\log_2(\frac{x}{2}))$}
\vspace{-20pt}
\label{lst:task3_t}
\begin{minted}[linenos, numbersep=10pt, firstnumber=1]{matlab}
% Adds values table
addpath('./task3_values.m');

function D = calcFinDiffTable(V)
    n = length(V(1, :));
    D = zeros(n, n+1);
    D(:, 1) = V(1, :)';
    D(:, 2) = V(2, :)';

    for j = 3:n+1
        for i = 1:n-j+2
            D(i, j) = (D(i+1, j-1) - D(i, j-1));
        end
    end
end

D = calcFinDiffTable(V);
disp(D);
\end{minted}
\end{listing}

Точка $x^{(1)} = 1.32$ находится ближе к началу таблице значений функции $f$, следовательно, для вычисления значения функции в ней стоит следует использовать многочлен Ньютона первого вида. На листинге \ref{lst:task3_1} предаствлена программа, вычисляющая значение функции в точке $x^{(1)} = 1.32$.

\begin{listing}[H]
\caption{Вычисление значения функции $f$ в точке $x^{(1)} = 1.32$ с помощью интерполяционного многочлена Ньютона первого вида}
\vspace{-20pt}
\label{lst:task3_1}
\begin{minted}[linenos, numbersep=10pt, firstnumber=1]{matlab}
% Adds values table and calculates finite difference table
addpath('./task3_t.m');

x_interp = 1.32;

n = length(D(1, :));
P = D(1, 2);
q = (x_interp - first) / h;
term = 1;

for i = 3:n
    term = term * (q - (i - 2) + 1);
    P = P + D(1, i) * term / factorial(i - 2);
end
disp(P);
\end{minted}
\end{listing}

Таким образом,

\[
P_7(1.32) = 2.2643 \approx 2.26.
\]

Точка  $x^{(2)} = 3.84$, напротив, находится ближе к концу таблицы заченений функции $f$, следовательно, для вычисления значения функции в ней следует использовать многочлен Ньютона второго вида. На листинге \ref{lst:task3_2} предаствлена программа, вычисляющая значение функции в точке $x^{(2)} = 3.84$.

\begin{listing}[H]
\caption{Вычисление значения функции $f$ в точке $x^{(2)} = 3.84$ с помощью интерполяционного многочлена Ньютона второго вида}
\vspace{-20pt}
\label{lst:task3_2}
\begin{minted}[linenos, numbersep=10pt, firstnumber=1]{matlab}
% Adds values table and calculates finite difference table
addpath('./task3_t.m');

x_interp = 3.84;

n = length(D(1, :));
m = length(D(:, 1));
P = D(m, 2);
q = (x_interp - last) / h;
term = 1;

for i = 3:n
    i_shifted = i - 2;
    term = term * (q + i_shifted - 1);
    P = P + D(m - i_shifted, i) * term / factorial(i_shifted);
end
disp(P);
\end{minted}
\end{listing}

Таким образом,

\[
P_7(3.84) = 0.267 \approx 0.27.
\]

Далее требуется найти значение функции $f$ в отчке $x = 2.2$ c помошью подходящего многочлена Гаусса, Стирлинга или Бесселя по таблице значений (табл. \ref{table:values_2}). Будем рассмвтривать рассматривать узлы $1.0, \dots, 3.5$.

Заметим, что точка $x = 2.2$ находится близко к середине~--- $2.25$~--- интервала $[a;a+h]$. $a = 2.0$~--- центральный узел, $h = 0.5$~--- шаг. Число узлов $n + 1 = 6$ (чётно), $n = 5$ (нечётно). Следовательно, для рассчётов можно использовать многочлен Бесселя.

Наглядно это показано на таблице \ref{table:bess}. Серые ячейки таблицы в вычислениях не используются. \fbox{Так} выделены необходимые для расчета разности.

\begin{table}[H]
    \centering
    \caption{Таблица конечных разностей (для рассчёта значения функции $f$ в точке $x = 2.2$ с помощью многочлена Бесселя)}
    \label{table:bess}
    \begin{tblr}{
        colspec = {X[c]X[c]X[c]X[c]X[c]X[c]X[c]X[c]},
        hlines,
        vlines,
        rows = {abovesep=5pt,belowsep=5pt, font=\small, $},
        row{1} = {font=\normal},
        cell{7}{1} = {yellow9},
        cell{2}{1} = {cmd=\underline},
        cell{4}{1} = {cmd=\underline},
        cell{6}{1} = {cmd=\underline},
        cell{8}{1} = {cmd=\underline},
        cell{10}{1} = {cmd=\underline},
        cell{12}{1} = {cmd=\underline},
        cell{6}{2} = {cmd=\fbox},
        cell{8}{2} = {cmd=\fbox},
        cell{7}{3} = {cmd=\fbox},
        cell{6}{4} = {cmd=\fbox},
        cell{8}{4} = {cmd=\fbox},
        cell{7}{5} = {cmd=\fbox},
        cell{6}{6} = {cmd=\fbox},
        cell{8}{6} = {cmd=\fbox},
        cell{7}{7} = {cmd=\fbox},
        column{8} = {bg={color=gray!50}},
        cell{8}{7} = {bg={color=gray!50}},
        cell{9}{6-7} = {bg={color=gray!50}},
        cell{10}{5-7} = {bg={color=gray!50}},
        cell{11}{4-7} = {bg={color=gray!50}},
        cell{12}{3-7} = {bg={color=gray!50}},
        row{13-14} = {bg={color=gray!50}}
    }
        x_i &
        y_i &
        \Delta y_i &
        \Delta^2 y_i &
        \Delta^3 y_i &
        \Delta^4 y_i &
        \Delta^5 y_i &
        \Delta^6 y_0 \\

        1.0000      &  3.1416 &         &         &         &         &         &        \\
                    &         & -1.1428 &         &         &         &         &        \\
        1.5000      &  1.9988 &         &  0.7148 &         &         &         &        \\
                    &         & -0.4280 &         & -0.6146 &         &         &        \\
        a = 2.0     &  1.5708 &         &  0.1002 &         &  0.5451 &         &        \\
        x = 2.2     &         & -0.3278 &         & -0.0695 &         & -0.5240 &        \\
        2.5000      &  1.2430 &         &  0.0307 &         &  0.0211 &         & 0.2528 \\
                    &         & -0.2971 &         & -0.0485 &         & -0.2712 &        \\
        3.0000      &  0.9460 &         & -0.0178 &         & -0.2501 &         &        \\
                    &         & -0.3148 &         & -0.2986 &         &         &        \\
        3.5000      &  0.6311 &         & -0.3163 &         &         &         &        \\
                    &         & -0.6311 &         &         &         &         &        \\
        4.0000      &       0 &         &         &         &         &         &        \\
    \end{tblr}
\end{table}

Интерполяционный многочлен Бесселя иммет вид

\begin{align*}
    P_n(a + qh) = \frac{y_0 + y_1}{2} + \frac{q - \frac{1}{2}}{1!}\Delta f_0 + \frac{q(q - 1)}{2!} \cdot \frac{\Delta^2 f_{-1} + \Delta^2 f_0}{2} + \\
    + \frac{(q - \frac{1}{2})q(q - 1)}{3!} \Delta^3f_{-1} + \dots + \\
    + \frac{q(q - 1)(q + 1)\dots(q + \frac{n-3}{2})(q - \frac{n-1}{2})}{(n-1)!} \cdot \frac{\Delta^{n-1}f_{-\frac{n-1}{2}} + \Delta^{n-1}f_{-\frac{n-3}{2}}}{2} + \\
    + \frac{(q - \frac{1}{2})q(q - 1)(q + 1)\dots(q + \frac{n - 3}{2}){q - \frac{n - 1}{2}}}{n!}\Delta^nf_{-\frac{n-1}{2}}. \\
    q = \frac{x - a}{h} = \frac{2.2 - 2.0}{0.5} = 0.4.
\end{align*}

На листинге \ref{lst:task3_3} приведена прогамма, вычисляющая значение многочлена Бесселя для заданных значений $x = 2.2, a = 2.0, h = 0.5$ и таблице конечных разностей \ref{table:bess}.

\begin{listing}[H]
\caption{программа, вычисляющая значение многочлена Бесселя}
\vspace{-20pt}
\label{lst:task3_3}
\begin{minted}[linenos, numbersep=10pt, firstnumber=1]{matlab}
% Adds values table and calculates finite difference table
addpath('./task3_t.m');

x_interp = 2.2;
a = 2.0;
n = 5;
q = (x_interp - a) / h;

P = (D(3, 2) + D(4, 2)) / 2 + (q - 0.5) / factorial(1) ...
    * D(3, 3) + q * (q - 1) / factorial(2) * (D(2, 4) + D(3, 4)) / 2 ...
    + (q - 1/2) * q * (q - 1) / factorial(3) * D(2, 5) ...
    + q * (q - 1) * (q + 1) * (q - 2) / factorial(4) ...
    * (D(1, 6) + D(2, 6)) / 2 ...
    + (q - 1/2) * q * (q - 1) * (q + 1) * (q - 2) / factorial(5) ...
    * D(1, 7);

disp(P);
\end{minted}
\end{listing}

Таким образом,

\[
P_6(2.2) = 1.4381 \approx 1.44.
\]

\section{Вывод}

В результате выполнения лабораторной работы были посчитаны прибоижённые значений функции, заданной по таблице, в точке с помощью интерполяционного многочлена Лагранжа и схемы Эйткена. Также были посчитаны аначения в точках с помощью первого и вторго многочленов нтютона Ньютона с разделёнными и конечными разностями. И было посчитано знаяение функции в точке с помощью многочлена Бесселя.

Также были освоены принципы работы и Matlab и базовые возможности языка; все функции написаны в нём (в приложении приведён исходный код всех программ).
